\section{Certificate conservatism and a broader-model precision boundary}
\label{sec:precision-boundary}

Two questions must be separated. A radius can fail to shrink because its
particular certificate is conservative, even at an all-valid distribution.
A different question is whether any uniformly valid procedure can improve
over the whole class of arbitrary bounded invalid means. The following
calculations address these questions separately. Neither establishes the
optimal adaptive length at an all-valid distribution.

\subsection{An explicit floor in the current certificate}

Consider $N$ observations per arm with an intercept-only balancing fit, so
$\gamma_i=1/N$, $h_i=1/N$ and $d_i=1/\{N(N-1)\}$. Under the bounded-invalid
certificate, $\rho=\sqrt N/2$ and the commutator term is zero. Thus
\[
 B_j=\frac1{4(N-1)}.
\]
Declare $g=(1-\varepsilon)K$ valid sources per arm, with
$0<\varepsilon<1/2$, although in this diagnostic all sources are actually
valid and have identical constant means. The top-$\varepsilon K$ certificate
gives
\[
 a_KC_B\longrightarrow\frac{\varepsilon}{4(N-1)}.
\]
The uncorrected empirical dispersion tends to zero, while the concentration
constants vanish as $K\to\infty$ at fixed $N$. Consequently the arm radius
converges to $\varepsilon/\{2\sqrt{(1-\varepsilon)(N-1)}\}$. Summing the two
arms yields
\[
 B_A\longrightarrow\frac{\varepsilon}{\sqrt{(1-\varepsilon)(N-1)}}.
\]
The contrast radius tends to
$\varepsilon/\sqrt{(1-2\varepsilon)(N-1)}$, which is larger. The current
adaptive certificate therefore retains the arm radius above. This proves a
floor for this \emph{source-band certificate}, including on favorable
all-valid data. The final target intersection can be shorter; the calculation
is not a lower bound for every inference method.

Additional structure can remove this particular cost. In the three-support-point
experiment of Section~\ref{sec:approximation-control}, the prespecified basis
$(1,x,x^2)$ spans every conditional mean on that support. It preserves uniform
balance weights and removes projection bias for all nonlinear amplitudes,
without supplying the true amplitude. A paired experiment compares this
saturated fit with bounded-invalid protection using $(1,x)$. The former
operates in a finite-support class with repeated observations at each point.
It does not solve arbitrary nonlinear means on distinct continuous designs.

\subsection{A conditional worst-case lower bound with a target anchor}

The contamination-neighborhood argument of \citet{chen2016} gives a useful
starting point. Its direct application to the present setting needs an
explicit embedding, because our patients are independent conditional on their
fixed site parameters. We provide that embedding here.

Fix distinct covariates within every treated arm, with $N_s$ treated patients
per source and $N_t$ at the target. They may be chosen to have full-rank,
low-leverage linear designs. Set all control means equal to $.4$ in both
worlds. In world $b\in\{0,1\}$, target and valid-source treated means are
constant $\theta_b$, where
\[
 \theta_0=\tfrac12,\qquad \theta_1=\tfrac12+\Delta,\qquad
 \Delta=\frac{c}{\sqrt{\max(N_s,N_t)}}.
\]
Thus the two target effects differ by $\Delta$. At an invalid source allow
an arbitrary vector of bounded treated means on its distinct covariates.
These vectors can be extended to bounded measurable functions, with no
smoothness requirement.

Let $P_b=\operatorname{Bernoulli}(\theta_b)^{\otimes N_s}$ and choose
$0<\varepsilon<1-\rho$, where the guaranteed valid count is
$g=\lceil\rho K\rceil$. If
\[
 t_s:=\operatorname{TV}(P_0,P_1)\le\frac{\varepsilon}{1-\varepsilon},
\]
there is a distribution $H$ and probability distributions $Q_b$ such that
\[
 H=(1-\varepsilon)P_0+\varepsilon Q_0
  =(1-\varepsilon)P_1+\varepsilon Q_1.
\]
For example, with probability mass functions $p_b$, take
\[
 h=(1-\varepsilon)\max(p_0,p_1)
       +\{1-(1-\varepsilon)(1+t_s)\}p_0,
 \qquad q_b=\{h-(1-\varepsilon)p_b\}/\varepsilon.
\]
All masses are nonnegative and sum to one. The source mixtures are identical
in the two worlds, however large $K$ becomes.

To realize $Q_b$ without introducing dependent patients in a fixed model,
draw a binary vector $y$ from $Q_b$ in the \emph{prior over site mean vectors}
and set that invalid site's patient means to $y$. Conditional on this mean
vector its Bernoulli outcomes are degenerate and independent. Mixing over
these vectors reproduces $Q_b$. This is a lower-bound prior, not the
outcome law of one fixed invalid site. Valid-site flags have probability
$1-\varepsilon$; condition each prior on at least $g$ such flags. Every
parameter configuration in the resulting prior satisfies the deterministic
valid-count assumption. If
\[
 q_K=\Prb\{\operatorname{Binomial}(K,1-\varepsilon)<g\},
\]
conditioning changes either source-data distribution by at most $q_K$ in
total variation. The two conditioned full-data mixtures, including the
independent target, consequently have distance at most
\[
 t_t+2q_K,\qquad
 t_t=\operatorname{TV}\{\operatorname{Bernoulli}(\theta_0)^{\otimes N_t},
                        \operatorname{Bernoulli}(\theta_1)^{\otimes N_t}\}.
\]

\begin{proposition}[Worst-case conditional length in the unrestricted invalid class]
An interval with conditional coverage at least $1-\alpha$ at every fixed
parameter configuration of the preceding class satisfies
\[
 \sup_P\E_P|I|\ \ge
 \Delta\,[1-2\alpha-t_t-2q_K]_+.
\]
For fixed $\varepsilon,\rho$ with $\rho<1-\varepsilon$, sufficiently small
constant $c>0$, and comparable $N_s,N_t$, the right-hand side has order
$N_s^{-1/2}$ as $K\to\infty$.
\end{proposition}

\begin{proof}
Coverage at every fixed configuration implies coverage under each conditioned
prior. Transfer the world-1 coverage event to world 0 using total variation.
Under world 0 both target effects are then in $I$ with probability at least
$1-2\alpha-t_t-2q_K$, forcing length at least $\Delta$. A prior-average length
cannot exceed the supremum over its fixed configurations. Finally,
$\operatorname{KL}\{\operatorname{Bernoulli}(1/2),
\operatorname{Bernoulli}(1/2+\Delta)\}
=-\tfrac12\log(1-4\Delta^2)$. Pinsker's inequality bounds both $t_s$ and
$t_t$ by a constant times $c$. Also
$q_K\le\exp\{-2K(1-\varepsilon-\rho)^2\}$, since fewer than
$\lceil\rho K\rceil$ valid flags implies their fraction is below $\rho$.
\end{proof}

This result says that a uniform fourth-root improvement over the entire
arbitrary-bounded-invalid \emph{conditional} class is impossible when target
and source arm sizes are comparable. The trusted target prevents an
identification failure, but still has its own sampling scale. If the target
is much larger, the displayed lower bound becomes correspondingly smaller.

The claim is deliberately limited: it is a worst-case conditional-design
result using unrestricted invalid mean functions. It is not an adaptive lower
bound at all-valid distributions, a lower bound for the weak-shift-only
subclass, or a proved random-i.i.d.-design population lower bound. In
particular, it does not show that the explicit certificate floor above is
optimal at favorable data. The code verifies the mixture identity and exact
binomial total variations, and reports the finite-$K$ count-conditioning cost.
